% ===== PRÉAMBULE CANONIQUE QCM — MATHÉMATIQUES =====
% Sert à compiler controle.tex ET à la revalidation automatique du JSON
% (valider_qcm.py). NE PAS MODIFIER : packages figés.
\documentclass[11pt,a4paper]{article}
\usepackage[utf8]{inputenc}\usepackage[T1]{fontenc}\usepackage{lmodern}
\usepackage[french]{babel}
\usepackage{amsmath,amssymb,mathtools}\usepackage{icomma}
\usepackage{siunitx}\sisetup{locale=FR}
\usepackage{xcolor}\usepackage{colortbl}
\usepackage{tikz}\usepackage{pgfplots}\pgfplotsset{compat=1.18}
\usepackage{enumitem}\usepackage{array,booktabs}
\usepackage[a4paper,margin=2.2cm]{geometry}
\usetikzlibrary{arrows.meta,calc,angles,quotes,positioning,patterns,backgrounds,shapes.geometric,intersections,babel}
\newcommand{\R}{\mathbb{R}}\newcommand{\N}{\mathbb{N}}\newcommand{\Z}{\mathbb{Z}}
\newcommand{\Q}{\mathbb{Q}}\newcommand{\C}{\mathbb{C}}\newcommand{\ds}{\displaystyle}
\newcommand{\vv}[1]{\vec{#1}}
% ===================================================
\begin{document}

\noindent\textbf{MATH-F05-Q01 --- Nombre dérivé : pente de la tangente}\par
Soit $f$ une fonction dérivable sur $\R$ telle que $f(1)=2$. La tangente à la courbe de $f$ au point d'abscisse $x_0=1$ passe également par le point de coordonnées $(4;4)$. Quel est le nombre dérivé $f'(1)$ ?
\begin{enumerate}[label=\Alph*)]
\item $\dfrac{3}{2}$
\item $\dfrac{4}{3}$
\item $\dfrac{2}{3}$
\item $\dfrac{1}{2}$
\end{enumerate}
\textit{Bonne réponse : C}\par
La tangente en $x_0=1$ a pour pente le nombre dérivé $f'(1)$. Elle passe par $\big(1;f(1)\big)=(1;2)$ et par $(4;4)$, donc $f'(1)=\dfrac{4-2}{4-1}=\dfrac{2}{3}$.\par
\textit{Pièges :}
\begin{itemize}
\item A : pente inversée $\dfrac{\Delta x}{\Delta y}=\dfrac{4-1}{4-2}=\dfrac{3}{2}$ (numérateur et dénominateur permutés).
\item B : oubli de soustraire $f(x_0)=2$, en calculant $\dfrac{4}{4-1}=\dfrac{4}{3}$.
\item D : oubli de soustraire $x_0=1$, en calculant $\dfrac{4-2}{4}=\dfrac{1}{2}$.
\end{itemize}
\vspace{8pt}\hrule\vspace{8pt}

\noindent\textbf{MATH-F05-Q02 --- Dérivée d'un quotient avec logarithme}\par
Soit $f$ la fonction définie sur $]0;+\infty[$ par $f(x)=\dfrac{\ln x}{x}$. Parmi les expressions suivantes, laquelle définit correctement la dérivée $f'(x)$ ?
\begin{enumerate}[label=\Alph*)]
\item $\dfrac{1-\ln x}{x^2}$
\item $\dfrac{\ln x-1}{x^2}$
\item $\dfrac{1-\ln x}{x}$
\item $\dfrac{1+\ln x}{x^2}$
\end{enumerate}
\textit{Bonne réponse : A}\par
Avec $U=\ln x$ ($U'=\tfrac1x$) et $V=x$ ($V'=1$) : $f'=\dfrac{U'V-UV'}{V^2}=\dfrac{\tfrac1x\cdot x-\ln x}{x^2}=\dfrac{1-\ln x}{x^2}$.\par
\textit{Pièges :}
\begin{itemize}
\item B : ordre du numérateur inversé $\dfrac{UV'-U'V}{V^2}=\dfrac{\ln x-1}{x^2}$.
\item C : dénominateur non élevé au carré, $\dfrac{1-\ln x}{x}$.
\item D : signe $+$ au lieu de $-$ au numérateur, $\dfrac{1+\ln x}{x^2}$.
\end{itemize}
\vspace{8pt}\hrule\vspace{8pt}

\noindent\textbf{MATH-F05-Q03 --- Dérivée composée d'une racine}\par
Soit $f$ la fonction définie sur $\R$ par $f(x)=\sqrt{x^2+1}$. Quelle est l'expression de $f'(x)$ ?
\begin{enumerate}[label=\Alph*)]
\item $\dfrac{2x}{\sqrt{x^2+1}}$
\item $\dfrac{x}{\sqrt{x^2+1}}$
\item $\dfrac{1}{2\sqrt{x^2+1}}$
\item $x\sqrt{x^2+1}$
\end{enumerate}
\textit{Bonne réponse : B}\par
En écrivant $f=(x^2+1)^{1/2}$, la règle $\big(U^n\big)'=n\,U^{n-1}U'$ avec $U=x^2+1$, $U'=2x$ et $n=\tfrac12$ donne $f'=\tfrac12(x^2+1)^{-1/2}\cdot 2x=\dfrac{x}{\sqrt{x^2+1}}$.\par
\textit{Pièges :}
\begin{itemize}
\item A : facteur $\tfrac12$ de la dérivée de la racine oublié, $\dfrac{2x}{\sqrt{x^2+1}}$.
\item C : facteur de chaîne $U'=2x$ oublié, $\dfrac{1}{2\sqrt{x^2+1}}$.
\item D : exposant non diminué de $1$, $\tfrac12(x^2+1)^{1/2}\cdot 2x=x\sqrt{x^2+1}$.
\end{itemize}
\vspace{8pt}\hrule\vspace{8pt}

\noindent\textbf{MATH-F05-Q04 --- Dérivée d'un produit avec exponentielle}\par
Soit $f$ la fonction définie sur $\R$ par $f(x)=x^2e^{-x}$. Quelle est l'expression de $f'(x)$ ?
\begin{enumerate}[label=\Alph*)]
\item $(2x-x^2)e^{-x}$
\item $(2x+x^2)e^{-x}$
\item $-x^2e^{-x}$
\item $2x\,e^{-x}$
\end{enumerate}
\textit{Bonne réponse : A}\par
Produit $UV$ avec $U=x^2$ ($U'=2x$) et $V=e^{-x}$ ($V'=-e^{-x}$) : $f'=2x\,e^{-x}+x^2(-e^{-x})=(2x-x^2)e^{-x}$.\par
\textit{Pièges :}
\begin{itemize}
\item B : signe de $(e^{-x})'$ pris positif, $(2x+x^2)e^{-x}$.
\item C : premier terme $U'V$ oublié, $-x^2e^{-x}$.
\item D : second terme $UV'$ oublié, $2x\,e^{-x}$.
\end{itemize}
\vspace{8pt}\hrule\vspace{8pt}

\noindent\textbf{MATH-F05-Q05 --- Dérivée trigonométrique composée}\par
Soit $f$ la fonction définie sur $\R$ par $f(x)=\cos(3x^2)$. Quelle est l'expression de $f'(x)$ ?
\begin{enumerate}[label=\Alph*)]
\item $6x\sin(3x^2)$
\item $-6x\sin(3x^2)$
\item $-6x\cos(3x^2)$
\item $-\sin(3x^2)$
\end{enumerate}
\textit{Bonne réponse : B}\par
$\big(\cos U\big)'=-U'\sin U$ avec $U=3x^2$ et $U'=6x$ : $f'(x)=-6x\sin(3x^2)$.\par
\textit{Pièges :}
\begin{itemize}
\item A : signe $-$ de la dérivée du cosinus oublié, $6x\sin(3x^2)$.
\item C : cosinus conservé au lieu de $-\sin$, $-6x\cos(3x^2)$.
\item D : facteur de chaîne $U'=6x$ oublié, $-\sin(3x^2)$.
\end{itemize}
\vspace{8pt}\hrule\vspace{8pt}

\noindent\textbf{MATH-F05-Q06 --- Sens de variation (signe de la dérivée)}\par
Soit $f$ la fonction définie sur $\R$ par $f(x)=x^3-12x$. Sur quel intervalle la fonction $f$ est-elle strictement décroissante ?
\begin{enumerate}[label=\Alph*)]
\item $]-\infty;-2[\,\cup\,]2;+\infty[$
\item $]-2;+\infty[$
\item $]-\infty;2[$
\item $]-2;2[$
\end{enumerate}
\textit{Bonne réponse : D}\par
$f'(x)=3x^2-12=3(x-2)(x+2)$ s'annule en $-2$ et $2$ ; le trinôme s'ouvrant vers le haut, $f'(x)<0$ (donc $f$ décroissante) exactement sur $]-2;2[$.\par
\textit{Pièges :}
\begin{itemize}
\item A : intervalle où $f'>0$ (croissance) au lieu de $f'<0$ (décroissance).
\item B : inéquation du second degré traitée comme $x>-2$ seulement.
\item C : inéquation du second degré traitée comme $x<2$ seulement.
\end{itemize}
\vspace{8pt}\hrule\vspace{8pt}

\noindent\textbf{MATH-F05-Q07 --- Extremums locaux (signe de la dérivée)}\par
Soit $f$ la fonction définie sur $\R$ par $f(x)=\dfrac{x^4}{4}-2x^2$. Combien d'extremums locaux la fonction $f$ possède-t-elle ?
\begin{enumerate}[label=\Alph*)]
\item $1$
\item $2$
\item $3$
\item $4$
\end{enumerate}
\textit{Bonne réponse : C}\par
$f'(x)=x^3-4x=x(x-2)(x+2)$ s'annule en $-2$, $0$ et $2$, en changeant de signe à chaque fois : $f$ possède 3 extremums locaux (un maximum en $0$, deux minimums en $-2$ et $2$).\par
\textit{Pièges :}
\begin{itemize}
\item A : une seule racine retenue ($x=2$), en perdant $x=0$ et $x=-2$.
\item B : $f'$ simplifiée par $x$, la racine $x=0$ est perdue, d'où 2 extremums.
\item D : confusion avec les 4 intervalles de monotonie.
\end{itemize}
\vspace{8pt}\hrule\vspace{8pt}

\noindent\textbf{MATH-F05-Q08 --- Dérivée seconde et point d'inflexion}\par
Soit $f$ la fonction définie sur $\R$ par $f(x)=x^3-6x^2+5$. En quelle abscisse la courbe de $f$ admet-elle un point d'inflexion ?
\begin{enumerate}[label=\Alph*)]
\item $x=0$
\item $x=2$
\item $x=4$
\item $x=1$
\end{enumerate}
\textit{Bonne réponse : B}\par
$f'(x)=3x^2-12x$ puis $f''(x)=6x-12$ s'annule en $x=2$ en changeant de signe : la courbe y admet un point d'inflexion.\par
\textit{Pièges :}
\begin{itemize}
\item A : $f'(x)=0$ résolue à la place de $f''$, racine $x=0$.
\item C : $f'(x)=0$ résolue à la place de $f''$, racine $x=4$.
\item D : dérivée seconde erronée $f''=6x-6$ (terme $-6x^2$ mal dérivé), $x=1$.
\end{itemize}
\vspace{8pt}\hrule\vspace{8pt}

\noindent\textbf{MATH-F05-Q09 --- Équation de la tangente}\par
Soit $f$ la fonction définie sur $\R$ par $f(x)=x^2-4x+1$. Quelle est l'équation de la tangente à la courbe de $f$ au point d'abscisse $x_0=3$ ?
\begin{enumerate}[label=\Alph*)]
\item $y=2x-2$
\item $y=6x-20$
\item $y=2x-4$
\item $y=2x-8$
\end{enumerate}
\textit{Bonne réponse : D}\par
$f(3)=9-12+1=-2$ et $f'(x)=2x-4$ donne $f'(3)=2$. Tangente : $y=(x-3)\cdot 2+(-2)=2x-8$.\par
\textit{Pièges :}
\begin{itemize}
\item A : décalage $(x-x_0)$ oublié, $y=f'(x_0)\,x+f(x_0)=2x-2$.
\item B : pente $f'(3)=6$ (terme $-4$ de $f'=2x-4$ oublié), $y=6x-20$.
\item C : erreur de signe $f(3)=+2$, $y=2(x-3)+2=2x-4$.
\end{itemize}
\vspace{8pt}\hrule\vspace{8pt}

\noindent\textbf{MATH-F05-Q10 --- Lecture graphique de la dérivée}\par
La courbe ci-dessous représente une fonction $f$ dérivable sur $\R$, admettant un maximum local en $x=-2$ et un minimum local en $x=1$.
\begin{center}
\begin{tikzpicture}
\begin{axis}[width=8.5cm,height=5.6cm,axis lines=middle,xlabel={$x$},ylabel={$y$},
   xtick={-2,1},ytick=\empty,domain=-3.3:2.4,samples=80,
   ymin=-2,ymax=4,clip=true,tick label style={font=\footnotesize}]
\addplot[blue,line width=1.2pt]{x^3/3 + x^2/2 - 2*x};
\end{axis}
\end{tikzpicture}
\end{center}
Parmi les affirmations suivantes, laquelle est correcte ?
\begin{enumerate}[label=\Alph*)]
\item $f'(-3)>0$
\item $f'(0)>0$
\item $f'(-2)<0$
\item $f'(1)<0$
\end{enumerate}
\textit{Bonne réponse : A}\par
Sur $]-\infty;-2[$ la courbe monte, donc $f'>0$ et en particulier $f'(-3)>0$. Aux extremums $x=-2$ et $x=1$, la tangente est horizontale donc $f'=0$ ; sur $]-2;1[$ la courbe descend, donc $f'(0)<0$.\par
\textit{Pièges :}
\begin{itemize}
\item B : $f'$ crue positive sur $]-2;1[$ alors que $f$ y décroît ($f'<0$).
\item C : dérivée crue négative au maximum, alors qu'un extremum donne $f'=0$.
\item D : dérivée crue négative au minimum, alors que $f'(1)=0$.
\end{itemize}
\vspace{8pt}\hrule\vspace{8pt}

\end{document}
